%
% File: abstract.tex
%
\chapter{Abstract}
\begin{SingleSpace}

This thesis presents the design, implementation, and operational preparation of an intelligent meal planning system for Kazakhstani households. The system generates weekly plans for households of one to ten members under hard budget, allergy, religious, time, and repetition constraints while jointly optimising preference, diversity, and multi-macronutrient fit. The core contribution is a three-tier hybrid planner: a deterministic candidate-pool builder enforces domain rules and quality gates; an OR-Tools CP-SAT model performs combinatorial assignment of recipes to (day, meal-slot) pairs; and a LightGBM listwise ranker, trained on real and synthetic feedback, replaces the deterministic scoring function when a trained artifact is available, with a graceful fallback chain otherwise. On the reported evaluation scenario, the planner reduces macro-target gaps from thousands of kilocalories and hundreds of grams of protein down to noise. The data layer is a reproducible ETL pipeline joining a Hugging Face recipe corpus, scraped supplementary sources, a United States Department of Agriculture (USDA) FoodData Central nutrition catalog, and a Kazakhstan-specific price baseline derived from official socially-significant-goods statistics and an Arbuz.kz retail snapshot; embedding-based clustering, MinHash deduplication, and an eight-component quality score yield a published catalog of more than ninety thousand recipes. The system is delivered as five sibling codebases coordinated by a single versioned OpenAPI contract: a Django REST backend with asynchronous Celery orchestration, a React Native mobile client built with the Expo managed workflow and structured around twenty-seven screens, a React + Vite web client for landing and marketing, a Terraform-managed infrastructure layer, and a separate ETL repository. The deployment layer is a Terraform-managed Amazon Web Services (AWS) infrastructure for three environments wired through GitHub OpenID Connect (OIDC), ECS Fargate, RDS Multi-AZ, ElastiCache, CloudFront, WAF, and Secrets Manager, accompanied by eight CI/CD workflows, five operational runbooks, and observability through Sentry, PostHog, and CloudWatch. Five Architecture Decision Records justify the lasting choices: monorepo, Fargate, hybrid CP-SAT plus learning-to-rank, Expo managed workflow, and PostgreSQL. The result is a single decision-support application in which classical constraint programming, modern learning-to-rank, and disciplined production engineering are combined end-to-end and reproducibly from the repository state.

\keywords{meal planning, recommender systems, constraint programming, OR-Tools CP-SAT, learning-to-rank, LightGBM, Django, React Native, Expo, AWS ECS Fargate, Terraform, reproducible bootstrap, household decision support, local cuisine, Kazakhstan.}

\end{SingleSpace}
\clearpage
